\subsection{Boolean Jacobians}
\label{sec:boolean-jacobians}

Let $\Two = (\{\bot, \top\}, \join, \bot, \meet, \top, \neg)$ be the two-element Boolean algebra. Throughout
this section we adopt the join-semilattice reading of $\SemiLat$ (\secref{cmon}).

\subsubsection{Overview}
\label{sec:boolean-jacobians:overview}

Consider the following subcategory $\FinVect_\Two$ of $\SemiLat$:

\begin{itemize}
\item Objects are finite Boolean vector spaces $\Two^n$, arising as iterated biproducts
(\secref{boolean-jacobians:vectors}).
\item Morphisms $f: \Two^m \to \Two^n$ are $n \times m$ Boolean matrices; composition is matrix multiplication
in the $(\join, \meet)$-semiring (\secref{boolean-jacobians:matrices}).
\item Every morphism has a \emph{transpose} $f^{\!\top}: \Two^n \to \Two^m$ (also $\join$-preserving) and a
\emph{right adjoint} $f_*: \Two^n \to \Two^m$ (meet-preserving, with $f \adj f_*$); these are De Morgan
duals of each other (\secref{boolean-jacobians:transpose}, \secref{boolean-jacobians:de-morgan}).
\item The transpose embeds the category contravariantly into $\HeytConj$ as conjugate pairs; the right
adjoint embeds it contravariantly into $\LatGal$ as Galois connections
(\secref{boolean-jacobians:embeddings}).
\item The conjugate pair $(f, f^{\!\top})$ is the Boolean analogue of the forward/reverse-mode Jacobian pair
in AD: the transpose $f^{\!\top}$ corresponds to the forward-mode Jacobian $\pushf{\phi}_x$ and the given $f$
plays the role of the reverse-mode $\pullf{\phi}_x$. The adjoint map $f_*$ seems to be a different analysis
with no direct real-AD analogue (\secref{boolean-jacobians:ad}).
\end{itemize}

\noindent The notation $\FinVect_\Two$ is a mild abuse, since $\Two$ is a semiring rather than a field.

\subsubsection{Boolean vectors as iterated biproducts}
\label{sec:boolean-jacobians:vectors}

Define $\Two^n$ inductively as the $n$-fold biproduct of $\Two$ in $\SemiLat$: $\Two^0$ is the zero object and
$\Two^{n+1} = \Two \biprod \Two^n$. Its carrier is the set of length-$n$ vectors of Boolean values, with
pointwise $\join$ and $\bot$. The carrier additionally inherits pointwise $\meet$ and $\top$ from $\Two$, so
each $\Two^n$ carries the richer structure of a Boolean algebra, not merely a join semilattice, although
$\SemiLat$ morphisms are only required to preserve joins.

\subsubsection{Morphisms as Boolean matrices}
\label{sec:boolean-jacobians:matrices}

Because $\SemiLat$ is $\CMon$-enriched with biproducts, homs between biproduct objects admit a canonical
\emph{matrix representation}: a morphism $f: \Two^m \to \Two^n$ is determined by the $n \times m$ family of
morphisms
\begin{salign*}
   M_f[i,j] = \biproj_i \comp f \comp \biinj_j : \Two \to \Two
\end{salign*}
\noindent where $\biinj_j: \Two \to \Two^m$ and $\biproj_i: \Two^n \to \Two$ are the biproduct
injections/projections. Since $\SemiLat(\Two, \Two)$ has exactly two morphisms (the identity and the zero
morphism), each entry $M_f[i,j]$ can be identified with an element of $\Two$, and $f$ corresponds to an $n
\times m$ matrix of $\Two$-values.

Equivalently, $f$ is determined by its values on the ``standard basis vectors'' $e_j = \biinj_j(\top):
\Two^m$, since $M_f[i,j]$ is precisely the $i$-th coordinate of $f(e_j)$. This reflects a more general fact:
for any $Y$ in $\SemiLat$, there is an isomorphism
\begin{salign*}
   \SemiLat(\Two^m, Y) \cong |Y|^m
\end{salign*}
\noindent (with $|Y|$ the carrier of $Y$), sending $f$ to $(f(e_1), \ldots, f(e_m))$. The isomorphism
follows from two steps: the universal property of the coproduct, $\SemiLat(\Two^m, Y) \cong
\prod_j \SemiLat(\Two, Y)$, and the observation that a $\join$-preserving map $\Two \to Y$ must send $\bot$
to $\bot$ and is otherwise determined by its value on $\top$, giving $\SemiLat(\Two, Y) \cong |Y|$ via
evaluation at $\top$. Concretely, writing $J(v) = \{j : v_j = \top\}$ for the set of indices at which $v$
is $\top$, we have
\begin{salign*}
   f(v) = \bigjoin_{j \in J(v)} f(e_j)
\end{salign*}
\noindent i.e.~the join of $f(e_j)$ over just those positions where $v$ is ``live''. This follows from $f$
being $\join$-preserving together with the biproduct decomposition $v = \bigjoin_j v_j \scalarmeet e_j$ (where
$v_j \scalarmeet e_j$ equals $e_j$ if $v_j = \top$ and $\bot$ otherwise). Specialising to $Y = \Two^n$ gives
$\SemiLat(\Two^m, \Two^n) \cong (\Two^n)^m$, i.e.\ $n \times m$ Boolean matrices, with composition given by
the usual matrix product in the $(\join, \meet)$-semiring.

We can think of $M_f$ as exactly the \emph{dependency relation} between the two index sets: $M_f[i,j] = \top$
iff position $i$ in the codomain is ``live'' in $f(e_j)$, i.e.\ iff position $i$ in the codomain depends on
position $j$ in the domain. Matrix composition, defined as usual in terms of the dot product (inner product)
in the $(\Two, \join, \meet)$-semiring:
\begin{salign*}
   (M_g \comp M_f)[i,k] = \bigjoin_j M_g[i,j] \meet M_f[j,k]
\end{salign*}
\noindent then corresponds to \emph{reachability}, as it evaluates to $\top$ iff there is some intermediate
position $j$ with both $M_g[i,j] = \top$ and $M_f[j,k] = \top$, i.e.~iff there is a chain $k \to j \to i$
witnessing that $i$ depends on $k$ via $j$.

\subsubsection{Transpose}
\label{sec:boolean-jacobians:transpose}

Since $\SemiLat$ is semi-additive, every morphism between biproducts has a \emph{transpose}. For
$f: \Two^m \to \Two^n$ with matrix $M_f$, the transpose $f^{\!\top}: \Two^n \to \Two^m$ has matrix
$(M_{f^{\!\top}})[j,i] = M_f[i,j]$. Concretely,
\begin{salign*}
   f^{\!\top}(w)_j = \bigjoin_{i} M_f[i,j] \meet w_i
\end{salign*}
\noindent The transpose is itself a morphism in $\SemiLat$: it is $\join$-preserving in its argument, just
like $f$.

\subsubsection{Two embeddings: $\HeytConj$ and $\LatGal$}
\label{sec:boolean-jacobians:embeddings}

We now treat $f: \Two^m \to \Two^n$ as the \emph{backward} direction of an analysis --- the map that
propagates output information back to input information (see \secref{boolean-jacobians:ad} for why this is
the appropriate reading). Accordingly, when we embed $f$ into $\HeytConj$ or $\LatGal$, the target
category's arrow points in the \emph{forward} direction $\Two^n \to \Two^m$, opposite to $f$. This choice
is what makes both embeddings \emph{contravariant} on morphisms: had we instead taken $f$ itself as the
forward direction, the embeddings would be covariant.

The two embeddings then pair $f$ with a forward-direction counterpart, in two different ways:

\begin{enumerate}
\item \emph{Conjugate pair.} The transpose $f^{\!\top}$ pairs with $f$ to form a conjugate pair in
$\HeytConj$ (\secref{categories-with-biproducts:heytconj}): both maps are $\join$-preserving, linked by
the disjointness condition
\begin{salign*}
   y \disj f^{\!\top}(x) \iff f(y) \disj x
\end{salign*}
\noindent Each $\Two^n$ is a Boolean, hence Heyting, algebra, so this assignment gives an embedding of a
subcategory of $\SemiLat$ into $\HeytConj$.

\item \emph{Galois connection.} The right adjoint $f_*: \Two^n \to \Two^m$ (meet-preserving) pairs with $f$
to form a Galois connection in $\LatGal$ (\secref{categories-with-biproducts:latgal}), with adjunction
\begin{salign*}
   y \le f_*(x) \iff f(y) \le x
\end{salign*}
\noindent Concretely, $f_*(x)_i = \bigmeet_j (M_f[i,j] \to x_j)$, where $a \to b = \neg a \join b$ is
Boolean implication. In the target arrow $\Two^n \to \Two^m$, $f$ is the left adjoint and $f_*$ the right
adjoint.
\end{enumerate}

\subsubsection{De Morgan duality and the derivation of the conjugate embedding}
\label{sec:boolean-jacobians:de-morgan}

The two backward maps are De Morgan duals:
\begin{salign*}
   \neg\,f^{\!\top}(x) = f_*(\neg x)
\end{salign*}
\noindent This is the Boolean-algebra identity relating the transpose (in the $(\join, \meet)$-semiring) to
the adjoint (in the dual $(\meet, \join)$-semiring), exploiting that $\Two^n$ is a Boolean algebra so that
$\neg$ is an involution.

Combined with the Boolean-algebra fact that $a \disj b \iff a \le \neg b$, we obtain the following chain:
\begin{salign*}
   y \disj f^{\!\top}(x)
   & \iff y \le \neg\,f^{\!\top}(x) \\
   & \iff y \le f_*(\neg x) \\
   & \iff f(y) \le \neg x \\
   & \iff f(y) \disj x
\end{salign*}
\noindent where the middle step is the Galois adjunction instantiated at $\neg x$. This exhibits the
conjugate embedding as a derived construction: the existence of a Galois connection $(f_*, f)$ together
with the Boolean structure on $\Two^n$ automatically yields the conjugate pair $(f^{\!\top}, f)$.

This is specific to the Boolean case. In a general Heyting algebra the equivalence $a \disj b \iff a \le
\neg b$ still holds (taking $\neg$ as the pseudo-complement), but the De Morgan identity $\neg f^{\!\top}
\simeq f_* \comp \neg$ relies on double-negation $\neg\neg = \id$, which is Boolean-specific. So in the
Boolean setting the conjugate and Galois perspectives on Jacobians are two faces of the same structure,
bridged by $\neg$.

\subsubsection{Connection to forward- and reverse-mode AD}
\label{sec:boolean-jacobians:ad}

Recall from \secref{auto-diff} that for a differentiable $\phi: \RR^p \to \RR^q$, forward mode applies the
Jacobian $\pushf{\phi}_x: \RR^p \linearto \RR^q$ to a tangent, and reverse mode applies its transpose
$\pullf{\phi}_x: \RR^q \linearto \RR^p$ to a cotangent. Both are $\RR$-linear maps in $\FinVect_\RR$.
Composition of derivatives is composition of linear maps in either direction, and the transpose is the
operation that flips the direction while preserving linearity: it is the dagger on the $\CMon$-enriched
biproduct category $\FinVect_\RR$.

The Boolean story is structurally identical at the level of linear algebra. The role of $\FinVect_\RR$ is
played by $\FinVect_\Two$, with biproducts $\Two^n$ and morphisms given by Boolean
matrices, and with the dagger given by matrix transposition. A morphism $g: \Two^p \to \Two^q$ is the
Boolean analogue of a linear map $\RR^p \to \RR^q$, and its transpose $g^{\!\top}: \Two^q \to \Two^p$ is the
Boolean analogue of the linear transpose; both are join-preserving in the same algebraic sense, just as
$\pushf{\phi}_x$ and $\pullf{\phi}_x$ are both $\RR$-linear.

Care is needed in lining up directions. In the preceding subsections, our given $f$ lives in
$\SemiLat$, and both embeddings of \secref{boolean-jacobians:embeddings} are \emph{contravariant} on
morphisms: $f: \Two^m \to \Two^n$ becomes part of an arrow $\Two^n \to \Two^m$ in $\HeytConj$ or $\LatGal$.
Reading the target category's arrow as the primal ``forward'' direction of a program, this means the
Jacobian proper --- the forward-mode linear map --- is not $f$ itself but its transpose
$f^{\!\top}: \Two^n \to \Two^m$: here $n$ indexes the inputs and $m$ indexes the outputs of that program. The given $f: \Two^m \to \Two^n$, going from outputs to inputs, is
the Boolean analogue of the \emph{reverse-mode} map $\pullf{\phi}_x$. Transposition is involutive, so this
is only a choice of labelling: the conjugate pair $(f, f^{\!\top})$ captures exactly the same content as
$(\pullf{\phi}_x, \pushf{\phi}_x)$, with the transpose flipping direction while preserving linearity.

The conjugate embedding is therefore the AD analogue in the standard sense: both sides use the dagger
(matrix transpose) to relate the two directions, and both directions live in the same linear category
($\FinVect_\RR$ or $\FinVect_\Two$). The Galois embedding is a genuinely different analysis. The right
adjoint $f_*: \Two^n \to \Two^m$ is not a morphism of $\FinVect_\Two$: it is $\meet$-preserving rather than
$\join$-preserving, and lives in the \emph{dual} semi-additive structure where $\meet$ plays the role of
addition. Concretely, reading $f_*$ in the standard sense gives the ``sufficiency'' interpretation (``is
every input that output $i$ depends on present?''), a universally quantified condition in a different
semiring.

More subtly, it is not just $f_*$ that is foreign to AD but the entire adjunction structure. The Galois
connection pairs $f$ with $f_*$ via $f \adj f_*$, and this adjunction has no analogue in $\FinVect_\RR$: the
reverse-mode map $\pullf{\phi}_x$ has no order-theoretic right adjoint, because $\RR$ is a field rather than
a lattice, and is related to $\pushf{\phi}_x$ only by transposition, not by adjunction. So while $f$ as a
linear map does correspond to $\pullf{\phi}_x$ under the dagger/conjugate analogy, $f$'s additional role as
the left adjoint of a Galois connection --- and the map $f_*$ that role brings with it --- is a
feature of the Boolean (or more generally, bounded lattice) setting with no AD counterpart.

The analogy is therefore:

\begin{center}
\begin{tabular}{ll}
 \textbf{Real AD} ($\FinVect_\RR$) & \textbf{Boolean AD} ($\FinVect_\Two$) \\
 Linear map $\pushf{\phi}_x: \RR^p \linearto \RR^q$ (forward mode) & $\join$-preserving $f^{\!\top}: \Two^n \to \Two^m$ \\
 Linear transpose $\pullf{\phi}_x: \RR^q \linearto \RR^p$ (reverse mode) & Given $f: \Two^m \to \Two^n$ \\
 Matrix multiplication (chain rule) & Composition in the $(\join, \meet)$-semiring \\
 Dagger on $\FinVect_\RR$ via inner product & Dagger on $\FinVect_\Two$ via $\bigjoin / \meet$ \\
\end{tabular}
\end{center}
